\begin{abstract}
Reconstructing and editing \yl{3D objects and scenes} both play crucial roles in computer graphics and computer vision.
Neural radiance fields \yl{(NeRFs) can achieve} %
realistic reconstruction and editing results but suffer from inefficiency in rendering. 
Gaussian splatting significantly accelerates rendering by rasterizing  Gaussian \yl{ellipsoids}.
However, \yl{Gaussian} splatting utilizes a single Spherical Harmonic (SH) function to model both texture and lighting,  limiting independent editing \yl{capabilities} of these components.  
Recently, attempts have been made to decouple texture and lighting with the \yl{Gaussian} splatting representation but may fail to produce plausible geometry and decomposition results on reflective scenes. 
Additionally, the \textit{forward shading} technique they employ introduces noticeable blending artifacts during relighting, as the geometry attributes of Gaussians are optimized under the original illumination and may not be suitable for novel lighting conditions. 
To address these issues, we introduce \name, a method for decoupling and editing the \yl{Gaussian} splatting representation using \textit{deferred shading}. 
To achieve successful decoupling, we model the illumination with a learnable environment map and define additional attributes such as texture parameters and normal direction on \yl{Gaussians}, where the normal is distilled from a jointly trained signed distance function. 
More importantly, we apply \textit{deferred shading}, resulting in more realistic relighting effects compared to previous methods. 
Both qualitative and quantitative experiments demonstrate the superior performance of \name in novel view synthesis and editing tasks. 
\end{abstract}
